\documentclass[10pt,a4paper]{amsart}
\usepackage{amsmath,amssymb,amsthm,graphicx}
\usepackage{verbatim}
\newenvironment{reference}{\comment}{\endcomment}
%\newenvironment{reference}{}{}
\newenvironment{slogan}{\comment}{\endcomment}
\newenvironment{history}{\comment}{\endcomment}

% For commutative diagrams we use Xy-pic
\usepackage[all]{xy}

% We use 2cell for 2-commutative diagrams.
\xyoption{2cell}
% Theorem environments.
%
\theoremstyle{plain}
\newtheorem{theorem}[subsection]{Theorem}
\newtheorem{proposition}[subsection]{Proposition}
\newtheorem{lemma}[subsection]{Lemma}

\theoremstyle{definition}
\newtheorem{definition}[subsection]{Definition}
\newtheorem{example}[subsection]{Example}
\newtheorem{exercise}[subsection]{Exercise}
\newtheorem{situation}[subsection]{Situation}

\theoremstyle{remark}
\newtheorem{remark}[subsection]{Remark}
\newtheorem{remarks}[subsection]{Remarks}

\numberwithin{equation}{subsection}

% Macros
%
\def\lim{\mathop{\mathrm{lim}}\nolimits}
\def\colim{\mathop{\mathrm{colim}}\nolimits}
\def\Spec{\mathop{\mathrm{Spec}}}
\def\Hom{\mathop{\mathrm{Hom}}\nolimits}
\def\Ext{\mathop{\mathrm{Ext}}\nolimits}
\def\SheafHom{\mathop{\mathcal{H}\!\mathit{om}}\nolimits}
\def\SheafExt{\mathop{\mathcal{E}\!\mathit{xt}}\nolimits}
\def\Sch{\mathit{Sch}}
\def\Mor{\mathop{\mathrm{Mor}}\nolimits}
\def\Ob{\mathop{\mathrm{Ob}}\nolimits}
\def\Sh{\mathop{\mathit{Sh}}\nolimits}
\def\NL{\mathop{N\!L}\nolimits}
\def\CH{\mathop{\mathrm{CH}}\nolimits}
\def\proetale{{pro\text{-}\acute{e}tale}}
\def\etale{{\acute{e}tale}}
\def\QCoh{\mathit{QCoh}}
\def\Ker{\mathop{\mathrm{Ker}}}
\def\Im{\mathop{\mathrm{Im}}}
\def\Coker{\mathop{\mathrm{Coker}}}
\def\Coim{\mathop{\mathrm{Coim}}}

% Boxtimes
%
\DeclareMathSymbol{\boxtimes}{\mathbin}{AMSa}{"02}

%
% Macros for moduli stacks/spaces
%
\def\QCohstack{\mathcal{QC}\!\mathit{oh}}
\def\Cohstack{\mathcal{C}\!\mathit{oh}}
\def\Spacesstack{\mathcal{S}\!\mathit{paces}}
\def\Quotfunctor{\mathrm{Quot}}
\def\Hilbfunctor{\mathrm{Hilb}}
\def\Curvesstack{\mathcal{C}\!\mathit{urves}}
\def\Polarizedstack{\mathcal{P}\!\mathit{olarized}}
\def\Complexesstack{\mathcal{C}\!\mathit{omplexes}}
% \Pic is the operator that assigns to X its picard group, usage \Pic(X)
% \Picardstack_{X/B} denotes the Picard stack of X over B
% \Picardfunctor_{X/B} denotes the Picard functor of X over B
\def\Pic{\mathop{\mathrm{Pic}}\nolimits}
\def\Picardstack{\mathcal{P}\!\mathit{ic}}
\def\Picardfunctor{\mathrm{Pic}}
\def\Deformationcategory{\mathcal{D}\!\mathit{ef}}
\usepackage{cleveref}
\crefname{originallemma}{original lemma}{original lemmas}
\Crefname{originallemma}{Original lemma}{Original lemmas}
\makeatletter
\newlabel{divisors-lemma-describe-reflexive-hull}{{0AY7}{source}{}{sourcepointer-1}{}}
\newlabel{divisors-lemma-describe-reflexive-hull@cref}{{[originallemma][0][]0AY7}{[1][0][]source}{}{}{}}
\makeatother
\crefname{originalreference}{original reference}{original references}
\Crefname{originalreference}{Original reference}{Original references}
\makeatletter
\newlabel{lemma-trace-higher}{{0AX5}{source}{}{sourcepointer-2}{}}
\newlabel{lemma-trace-higher@cref}{{[originalreference][0][]0AX5}{[1][0][]source}{}{}{}}
\makeatother
\title{Extension of the trace map on differential forms to the reflexive hull}
\author{The Stacks Project Authors}
\date{stacks-project:resolve.tex:local-snapshot-20261012; GNU FDL 1.2 or later}
\begin{document}
\maketitle
\noindent\textit{Editorial note.} All delivered mathematical text is the author's own, copied from the pinned Stacks Project chapter file resolve.tex; no statement or proof step was rewritten, reordered or removed.\par\medskip
\noindent\textit{Editorial note.} The counted proof is the author's complete argument: the coherence reduction to points of dimension one, the classification of the local purely inseparable extension of discrete valuation rings, and the application of the local trace construction.\par\medskip
\noindent\textit{Editorial note.} The label lemma-trace-extends names the counted statement, and the proof environment immediately following it in the source is delivered verbatim as its proof.\par\medskip
\noindent\textit{Editorial note.} The counted statement is carried by the author's own lemma environment in the chapter source.\par\medskip
\noindent\textit{Editorial note.} Pre-open gate: python operations/agent-authoring/tools/precard\_gate.py --source resolve.tex --label lemma-trace-extends --start 590 --end 636 --local-id 28100 returned RESULT=OK (L28100 free, occupancy FREE, public interval clear). Pre-assembly tag self-check: canonical tag resolved and not present in reference/public-items.jsonl.\par\medskip
\noindent\textit{Editorial note.} The statement and proof are copied from the pinned chapter file; the cited inputs are kept as external pointers with their canonical Stacks tags.\par\medskip
\noindent\textit{Allowed standard prerequisites (editorial).}\begin{itemize}
\item The trace map on differential forms for a degree p purely inseparable extension of rings. Same chapter, the lemma constructing the trace map on differential forms for a ring extension of the form a polynomial ring modulo a p-th power relation. Provides the trace map in the local model, which is the one that has to be extended to all of the source.
\item Coherence of the differentials and of their pushforward. The exact sequence relating the pullback of the differentials of the base, the differentials of the source, and the relative differentials. Shows that the modules involved are coherent, so that the extension problem can be checked at points of dimension one.
\item Extending a map on the reflexive hull by checking points of dimension one. Divisors, the lemma describing the reflexive hull in terms of localizations at points of dimension one. Reduces the construction of the trace map to the local case of a discrete valuation ring.
\item The structure of a finite purely inseparable extension of discrete valuation rings of degree p. The classification of such extensions into the ramified case and the unramified case with purely inseparable residue field extension, giving the presentation as a ring modulo a p-th power relation. Puts the local situation into the shape required by the trace construction.
\end{itemize}\medskip
\ifdefined\hypertarget\hypertarget{sourcepointer-1}{}\fi
\noindent\textit{Source reference.} Original lemma 0AY7: sources/stacks-chow/tags.txt, source lines 13780--13780. Prerequisite pointer invoked at proof line(s) 612; it names a standard fact of a sibling chapter that the counted argument uses, and every step of the counted proof built on it is retained verbatim in the delivered body.\par\medskip
\ifdefined\hypertarget\hypertarget{sourcepointer-2}{}\fi
\noindent\textit{Source reference.} Original reference 0AX5: sources/stacks-resolve/pinned-original-source.tex, source lines 501--501. Prerequisite pointer invoked at proof line(s) 604; it names a standard fact of the same chapter that the counted argument uses, and every step of the counted proof built on it is retained verbatim in the delivered body.\par\medskip
\begin{lemma}
\label{lemma-trace-extends}
Let $S$ be a scheme over $\mathbf{F}_p$. Let $f : Y \to X$ be a finite morphism
of Noetherian normal integral schemes over $S$. Assume
\begin{enumerate}
\item the extension of function fields is purely inseparable of degree $p$, and
\item $\Omega_{X/S}$ is a coherent $\mathcal{O}_X$-module (for example
if $X$ is of finite type over $S$).
\end{enumerate}
For $i \geq 1$ there is a canonical map
$$
\text{Tr} : f_*\Omega^i_{Y/S} \longrightarrow (\Omega_{X/S}^i)^{**}
$$
whose stalk in the generic point of $X$ recovers the trace map of
Lemma \ref{lemma-trace-higher}.
\end{lemma}

\begin{proof}
The exact sequence $f^*\Omega_{X/S} \to \Omega_{Y/S} \to \Omega_{Y/X} \to 0$
shows that $\Omega_{Y/S}$ and hence $f_*\Omega_{Y/S}$ are coherent modules
as well. Thus it suffices to prove the trace map in the generic point
extends to stalks at $x \in X$ with $\dim(\mathcal{O}_{X, x}) = 1$, see
Divisors, Lemma \ref{divisors-lemma-describe-reflexive-hull}.
Thus we reduce to the case discussed in the next paragraph.

\medskip\noindent
Assume $X = \Spec(A)$ and $Y = \Spec(B)$ with $A$ a discrete valuation
ring and $B$ finite over $A$. Since the induced extension $L/K$
of fraction fields is purely inseparable, we see that $B$ is local too.
Hence $B$ is a discrete valuation ring too. Then either
\begin{enumerate}
\item $B/A$ has ramification index $p$ and hence $B = A[x]/(x^p - a)$
where $a \in A$ is a uniformizer, or
\item $\mathfrak m_B = \mathfrak m_A B$ and the residue field
$B/\mathfrak m_A B$ is purely inseparable of degree $p$ over
$\kappa_A = A/\mathfrak m_A$.
Choose any $x \in B$ whose residue class is not in $\kappa_A$
and then we'll have $B = A[x]/(x^p - a)$ where $a \in A$ is
a unit.
\end{enumerate}
Let $\Spec(\Lambda) \subset S$ be an affine open such that
$X$ maps into $\Spec(\Lambda)$. Then we can apply
Lemma \ref{lemma-trace-higher}
to see that the trace map extends to
$\Omega^i_{B/\Lambda} \to \Omega^i_{A/\Lambda}$
for all $i \geq 1$.
\end{proof}

\end{document}
